\documentclass[12pt,letterpaper]{article}

% ---------- Paquetes ----------
\usepackage[utf8]{inputenc}
\usepackage[T1]{fontenc}
\usepackage[spanish,es-tabla]{babel}
\usepackage[margin=2.5cm]{geometry}
\usepackage{graphicx}
\usepackage{float}
\usepackage{amsmath,amssymb}
\usepackage{booktabs}
\usepackage{xcolor}
\usepackage{listings}
\usepackage{caption}
\usepackage{fancyhdr}
\usepackage[hidelinks]{hyperref}

\graphicspath{{imagenes/}}

% ---------- Datos de la práctica ----------
\newcommand{\numpractica}{XX}
\newcommand{\titulopractica}{Título de la práctica}
\newcommand{\fechaentrega}{\today}
\newcommand{\autor}{Domínguez Palacios Jesús Alejandro}
\newcommand{\profesor}{Nombre del profesor}
\newcommand{\grupo}{XX}

% ---------- Estilo de código ----------
\lstset{
  basicstyle=\ttfamily\small,
  keywordstyle=\color{blue},
  commentstyle=\color{gray},
  stringstyle=\color{red!60!black},
  numbers=left, numberstyle=\tiny, stepnumber=1,
  frame=single, breaklines=true, tabsize=4,
  showstringspaces=false
}

% ---------- Encabezado ----------
\pagestyle{fancy}
\fancyhf{}
\lhead{OAC -- Laboratorio}
\rhead{Práctica \numpractica}
\cfoot{\thepage}

\begin{document}

% ---------- Portada ----------
\begin{titlepage}
  \centering
  {\Large \textbf{Universidad Nacional Autónoma de México}\par}
  \vspace{0.3cm}
  {\large Facultad de Ingeniería\par}
  \vspace{1.5cm}
  {\large Laboratorio de Organización y Arquitectura de Computadoras\par}
  \vspace{2cm}
  {\LARGE \textbf{Práctica \numpractica}\par}
  \vspace{0.5cm}
  {\Large \titulopractica\par}
  \vfill
  \begin{tabular}{rl}
    \textbf{Alumno:}   & \autor \\
    \textbf{Profesor:} & \profesor \\
    \textbf{Grupo:}    & \grupo \\
    \textbf{Fecha:}    & \fechaentrega \\
  \end{tabular}
\end{titlepage}

\tableofcontents
\newpage

\section{Objetivo}
Escribe aquí el objetivo de la práctica.

\section{Introducción}
Marco teórico breve.

\section{Desarrollo}
Describe el procedimiento.

% Ejemplo de imagen:
% \begin{figure}[H]
%   \centering
%   \includegraphics[width=0.7\textwidth]{diagrama.png}
%   \caption{Descripción de la figura.}
% \end{figure}

% Ejemplo de código:
% \begin{lstlisting}[language=VHDL]
% -- código aquí
% \end{lstlisting}

\section{Resultados}
Muestra y analiza los resultados.

\section{Conclusiones}
Conclusiones de la práctica.

\begin{thebibliography}{9}
  \bibitem{ref1} Autor. \textit{Título}. Editorial, año.
\end{thebibliography}

\end{document}
